% =====================================================================
\chapter{Descenso de gradiente}
\label{tema:2}
% =====================================================================

El descenso de gradiente es el algoritmo del que derivan todos los optimizadores del módulo. Este tema demuestra sus garantías de convergencia en tres situaciones de generalidad decreciente: funciones suaves arbitrarias, funciones fuertemente convexas y funciones cuadráticas, en las que el análisis es exacto. El resultado central es que el número de iteraciones crece proporcionalmente al número de condición $\kappa$. Se sigue~\cite[§9.2 y §9.3]{bv}, \cite[§12.3]{d2l} y las clases~21 y~22 de~\cite{mit18065}.

\section{El método y el lema de descenso}

\begin{definicion}[Descenso de gradiente]\label{def:gd}
Dados $\theta_0\in\R^{n}$ y una \emph{tasa de aprendizaje} (o \emph{paso}) $\eta>0$, el \emph{descenso de gradiente} genera la sucesión
\[
\theta_{k+1}=\theta_k-\eta\,\nabla f(\theta_k),\qquad k=0,1,2,\dots
\]
\end{definicion}

La dirección $-\nabla f(\theta_k)$ es la de máximo descenso local (módulo~4, tema~1). Por el desarrollo de Taylor de primer orden, $f(\theta_k-\eta\nabla f(\theta_k))=f(\theta_k)-\eta\norm{\nabla f(\theta_k)}^{2}+o(\eta)$, de modo que un paso suficientemente pequeño reduce $f$ siempre que el gradiente no sea nulo. El lema siguiente cuantifica «suficientemente pequeño».

\begin{lema}[Lema de descenso]\label{lem:descenso}
Si $f$ es $L$-suave, para todo $x,y\in\R^{n}$,
\[
f(y)\le f(x)+\nabla f(x)\T(y-x)+\frac{L}{2}\norm{y-x}^{2}.
\]
\end{lema}

\begin{proof}
Sea $d=y-x$ y $\varphi(t)=f(x+td)$, con $\varphi'(t)=\nabla f(x+td)\T d$. Por el teorema fundamental del cálculo,
\[
f(y)-f(x)-\nabla f(x)\T d=\int_0^1\bigl(\nabla f(x+td)-\nabla f(x)\bigr)\T d\,dt.
\]
Por la desigualdad de Cauchy--Schwarz y la \cref{def:suave}, el integrando está acotado por $\norm{\nabla f(x+td)-\nabla f(x)}\,\norm{d}\le Lt\norm{d}^{2}$. Integrando, $\int_0^1 Lt\norm{d}^{2}\,dt=\frac{L}{2}\norm{d}^{2}$.
\end{proof}

\begin{corolario}[Descenso suficiente]\label{cor:descenso}
Si $f$ es $L$-suave y $\theta^{+}=\theta-\eta\nabla f(\theta)$ con $0<\eta<2/L$, entonces
\[
f(\theta^{+})\le f(\theta)-\eta\Bigl(1-\frac{L\eta}{2}\Bigr)\norm{\nabla f(\theta)}^{2}.
\]
Con $\eta=1/L$ la cota es $f(\theta^{+})\le f(\theta)-\frac{1}{2L}\norm{\nabla f(\theta)}^{2}$, y es la mejor que da el lema.
\end{corolario}

\begin{proof}
Basta tomar $x=\theta$ e $y=\theta^{+}$ en el \cref{lem:descenso}, con $y-x=-\eta\nabla f(\theta)$: el término lineal es $-\eta\norm{\nabla f}^{2}$ y el cuadrático, $\frac{L}{2}\eta^{2}\norm{\nabla f}^{2}$. La función $\eta\mapsto\eta(1-L\eta/2)$ es positiva en $(0,2/L)$ y alcanza su máximo, $1/(2L)$, en $\eta=1/L$.
\end{proof}

\begin{teorema}[Caso no convexo]\label{teo:gd-noconvexo}
Sea $f$ $L$-suave y acotada inferiormente por $f^{\ast}$. El descenso de gradiente con $\eta=1/L$ cumple, para todo $K\ge 1$,
\[
\min_{0\le k<K}\norm{\nabla f(\theta_k)}^{2}\le\frac{2L\bigl(f(\theta_0)-f^{\ast}\bigr)}{K}.
\]
\end{teorema}

\begin{proof}
Por el \cref{cor:descenso}, $\frac{1}{2L}\norm{\nabla f(\theta_k)}^{2}\le f(\theta_k)-f(\theta_{k+1})$. Sumando para $k=0,\dots,K-1$, la suma de la derecha es telescópica y vale $f(\theta_0)-f(\theta_K)\le f(\theta_0)-f^{\ast}$. El mínimo de $K$ números no supera su media, luego $\min_k\norm{\nabla f(\theta_k)}^{2}\le\frac1K\sum_k\norm{\nabla f(\theta_k)}^{2}\le 2L(f(\theta_0)-f^{\ast})/K$.
\end{proof}

El \cref{teo:gd-noconvexo} solo garantiza aproximarse a un punto crítico, que puede ser una silla; es la clase de garantía que cabe esperar para la pérdida de una red profunda.

\section{Convergencia lineal en funciones fuertemente convexas}

\begin{teorema}[Convergencia lineal]\label{teo:gd-fuerte}
Sea $f$ $\mu$-fuertemente convexa y $L$-suave, con valor óptimo $f^{\ast}$. El descenso de gradiente con $\eta=1/L$ cumple
\[
f(\theta_k)-f^{\ast}\le\Bigl(1-\frac{1}{\kappa}\Bigr)^{k}\bigl(f(\theta_0)-f^{\ast}\bigr),\qquad\kappa=\frac{L}{\mu}.
\]
\end{teorema}

\begin{proof}
Por el \cref{cor:descenso} y la desigualdad de Polyak--Łojasiewicz (\cref{prop:pl}),
\[
f(\theta_{k+1})-f^{\ast}\le f(\theta_k)-f^{\ast}-\frac{1}{2L}\norm{\nabla f(\theta_k)}^{2}\le f(\theta_k)-f^{\ast}-\frac{2\mu}{2L}\bigl(f(\theta_k)-f^{\ast}\bigr)=\Bigl(1-\frac{\mu}{L}\Bigr)\bigl(f(\theta_k)-f^{\ast}\bigr).
\]
Iterando $k$ veces se obtiene el enunciado.
\end{proof}

\begin{corolario}[Número de iteraciones]\label{cor:iteraciones}
En las condiciones del \cref{teo:gd-fuerte}, para alcanzar $f(\theta_k)-f^{\ast}\le\epsilon\,(f(\theta_0)-f^{\ast})$ basta con $k\ge\kappa\ln(1/\epsilon)$ iteraciones.
\end{corolario}

\begin{proof}
Como $\ln(1-x)\le -x$ para $x\in[0,1)$, se tiene $(1-1/\kappa)^{k}\le e^{-k/\kappa}$, que es $\le\epsilon$ si $k\ge\kappa\ln(1/\epsilon)$.
\end{proof}

La convergencia se llama \emph{lineal} porque el error se multiplica en cada paso por un factor constante menor que~1: el logaritmo del error decrece linealmente con $k$. El factor depende solo de $\kappa$. El análisis de~\cite[§9.3]{bv} obtiene el mismo factor $1-\mu/L$, salvo constantes, para búsquedas lineales exacta y con retroceso.

\section{Análisis exacto en cuadráticas}

En una función cuadrática el método es lineal en $\theta$ y puede resolverse en forma cerrada. Toda función de clase $C^{2}$ es aproximadamente cuadrática cerca de un mínimo no degenerado, de modo que este caso describe el comportamiento asintótico general.

\begin{ejemplo}[Una variable]\label{ej:gd-1d}
Sea $f(x)=\frac{a}{2}x^{2}$ con $a>0$. Entonces $x_{k+1}=x_k-\eta a x_k=(1-\eta a)\,x_k$, de donde $x_k=(1-\eta a)^{k}x_0$. El método converge si y solo si $\abs{1-\eta a}<1$, es decir, $0<\eta<2/a$. Para $f(x)=2x^{2}$ ($a=4$), $x_0=1$ y $\eta=0{,}1$, el factor es $0{,}6$ y los iterados son $1$; $0{,}6$; $0{,}36$; $0{,}216$; $0{,}1296$. Con $\eta=0{,}6>2/a=0{,}5$, el factor es $-1{,}4$: el iterado cambia de signo en cada paso y su módulo crece.
\end{ejemplo}

\begin{teorema}[Descenso de gradiente en cuadráticas]\label{teo:gd-cuadratica}
Sea $f(\theta)=\tfrac12\theta\T Q\theta$ con $Q$ simétrica de autovalores $0<\mu=\lambda_1\le\dots\le\lambda_n=L$. Sea $\rho(\eta)=\max_i\abs{1-\eta\lambda_i}$.
\begin{enumerate}
  \item $\norm{\theta_k}\le\rho(\eta)^{k}\norm{\theta_0}$, con igualdad si $\theta_0$ es autovector del autovalor que realiza el máximo.
  \item El método converge a $\theta^{\ast}=0$ para todo $\theta_0$ si y solo si $0<\eta<2/L$.
  \item El valor de $\eta$ que minimiza $\rho(\eta)$ es $\eta^{\ast}=2/(\mu+L)$, y
  \[
  \rho(\eta^{\ast})=\frac{L-\mu}{L+\mu}=\frac{\kappa-1}{\kappa+1}.
  \]
\end{enumerate}
\end{teorema}

\begin{proof}
Como $\nabla f(\theta)=Q\theta$, la iteración es $\theta_{k+1}=(I-\eta Q)\theta_k$, luego $\theta_k=(I-\eta Q)^{k}\theta_0$. Por el teorema espectral (módulo~3), $Q=U\Lambda U\T$ con $U$ ortogonal. En las coordenadas $z=U\T\theta$, la iteración se desacopla: $z^{(i)}_{k}=(1-\eta\lambda_i)^{k}z^{(i)}_{0}$.

(i) Como $U$ conserva la norma, $\norm{\theta_k}^{2}=\sum_i(1-\eta\lambda_i)^{2k}(z^{(i)}_0)^{2}\le\rho(\eta)^{2k}\norm{\theta_0}^{2}$, con igualdad en el caso indicado.

(ii) Por (i), la condición $\rho(\eta)<1$ es suficiente. Es necesaria: si $\abs{1-\eta\lambda_i}\ge 1$ para algún $i$, el autovector correspondiente no converge a~0. Y $\abs{1-\eta\lambda_i}<1$ para todo $i$ equivale a $0<\eta\lambda_i<2$ para todo $i$, es decir, a $0<\eta<2/L$.

(iii) Para $\eta$ fijo, $\lambda\mapsto\abs{1-\eta\lambda}$ es convexa, así que su máximo en $[\mu,L]$ se alcanza en un extremo: $\rho(\eta)=\max\{\abs{1-\eta\mu},\abs{1-\eta L}\}$. La primera función es decreciente en $\eta$ mientras $\eta\le 1/\mu$, y la segunda es creciente para $\eta\ge 1/L$; el máximo de ambas es mínimo donde se cortan, $1-\eta\mu=\eta L-1$, es decir, $\eta^{\ast}=2/(\mu+L)$. Sustituyendo, $\rho(\eta^{\ast})=1-2\mu/(\mu+L)=(L-\mu)/(L+\mu)$, y dividiendo por $\mu$ se obtiene $(\kappa-1)/(\kappa+1)$.
\end{proof}

\begin{observacion}
El \cref{teo:gd-cuadratica} explica la conducta del método. La dirección de curvatura máxima, $\lambda_n=L$, limita el paso: $\eta<2/L$. La dirección de curvatura mínima, $\lambda_1=\mu$, se contrae entonces con un factor $1-\eta\mu>1-2/\kappa$, próximo a~1 si $\kappa$ es grande. Con el paso óptimo, las dos direcciones extremas se contraen al mismo ritmo, pero con signos opuestos: la de curvatura alta oscila.
\end{observacion}

\begin{ejemplo}[Dos variables]\label{ej:gd-2d}
Sea $f(x,y)=x^{2}+5y^{2}$, con $Q=\diag(2,10)$, $\mu=2$, $L=10$ y $\kappa=5$. Con $\eta=0{,}05$, los factores son $1-0{,}1=0{,}9$ en $x$ y $1-0{,}5=0{,}5$ en $y$. Desde $(1,1)$:
\begin{center}\small
\begin{tabular}{@{}cccc@{}}
\toprule
$k$ & $x_k$ & $y_k$ & $f(x_k,y_k)$ \\
\midrule
0 & 1 & 1 & 6 \\
1 & 0{,}9 & 0{,}5 & 2{,}06 \\
2 & 0{,}81 & 0{,}25 & 0{,}9686 \\
3 & 0{,}729 & 0{,}125 & 0{,}6096 \\
\bottomrule
\end{tabular}
\end{center}
La coordenada $y$, de curvatura alta, converge deprisa; la $x$ se arrastra. El paso óptimo es $\eta^{\ast}=2/12=1/6$, con $\rho=4/6=2/3$; la \cref{fig:zigzag} muestra la trayectoria resultante.
\end{ejemplo}

\begin{figure}[h]
\centering
\begin{tikzpicture}[scale=2.6,font=\small]
  \foreach \c in {0.15,0.45,1,2,3.4,6}{
    \draw[regla,thin] (0,0) ellipse ({sqrt(\c)} and {sqrt(\c/5)});
  }
  \draw[->,gray] (-2.6,0) -- (2.6,0) node[right,black] {$x$};
  \draw[->,gray] (0,-1.2) -- (0,1.25) node[above,black] {$y$};
  \draw[thick,-{Stealth[length=2mm]}] (1,1) -- (0.6667,-0.6667);
  \draw[thick,-{Stealth[length=2mm]}] (0.6667,-0.6667) -- (0.4444,0.4444);
  \draw[thick,-{Stealth[length=2mm]}] (0.4444,0.4444) -- (0.2963,-0.2963);
  \draw[thick,-{Stealth[length=2mm]}] (0.2963,-0.2963) -- (0.1975,0.1975);
  \draw[thick,-{Stealth[length=2mm]}] (0.1975,0.1975) -- (0.1317,-0.1317);
  \foreach \p in {(1,1),(0.6667,-0.6667),(0.4444,0.4444),(0.2963,-0.2963),(0.1975,0.1975),(0.1317,-0.1317)}{\fill \p circle (0.6pt);}
  \node[above right] at (1,1) {$\theta_0$};
  \node[below right] at (0.6667,-0.6667) {$\theta_1$};
  \node[above right] at (0.4444,0.4444) {$\theta_2$};
  \fill (0,0) circle (0.8pt) node[below left] {$\theta^{\ast}$};
\end{tikzpicture}
\caption{Descenso de gradiente sobre $f(x,y)=x^{2}+5y^{2}$ con el paso óptimo $\eta^{\ast}=1/6$, desde $(1,1)$. Cada coordenada se multiplica por $2/3$ (en $x$) o por $-2/3$ (en $y$): la trayectoria zigzaguea entre las paredes del valle. Las elipses son curvas de nivel.}
\label{fig:zigzag}
\end{figure}

\begin{corolario}\label{cor:iter-cuadratica}
En las condiciones del \cref{teo:gd-cuadratica}, con el paso óptimo, el menor $k$ con $\rho^{k}\le\epsilon$ es $k=\lceil\ln\epsilon/\ln\rho\rceil$, que para $\kappa$ grande es aproximadamente $\frac{\kappa}{2}\ln(1/\epsilon)$.
\end{corolario}

\begin{proof}
La primera afirmación se obtiene tomando logaritmos en $\rho^{k}\le\epsilon$, teniendo en cuenta que $\ln\rho<0$. Para la segunda, $\ln\rho=\ln(1-1/\kappa)-\ln(1+1/\kappa)=-2/\kappa+O(\kappa^{-3})$.
\end{proof}

El \cref{tab:kappa} aplica el corolario con $\epsilon=10^{-3}$. La última columna anticipa el resultado del \cref{tema:4}: con momentum, el número de iteraciones crece como $\sqrt\kappa$ en lugar de como $\kappa$.

\begin{table}[h]
\centering\small
\caption{Factor de contracción e iteraciones necesarias para reducir el error a $10^{-3}$ en una cuadrática, con parámetros óptimos.}\label{tab:kappa}
\begin{tabular}{@{}rcrcr@{}}
\toprule
$\kappa$ & $\rho$ del gradiente & Iteraciones & $\rho$ con momentum (\cref{teo:heavy-ball}) & Iteraciones \\
\midrule
10 & $0{,}8182$ & 35 & $0{,}5195$ & 11 \\
100 & $0{,}9802$ & 346 & $0{,}8182$ & 35 \\
1000 & $0{,}9980$ & 3454 & $0{,}9387$ & 110 \\
\bottomrule
\end{tabular}
\end{table}

\section{La función de Rosenbrock y el borde de estabilidad}

\begin{ejemplo}[Función de Rosenbrock]\label{ej:rosenbrock}
La función $f(x,y)=(1-x)^{2}+100(y-x^{2})^{2}$ es no convexa y tiene su único mínimo en $(1,1)$, con $f=0$. Su hessiano en el mínimo es
\[
\nabla^{2}f(1,1)=\begin{pmatrix}802&-400\\-400&200\end{pmatrix},
\]
de autovalores $\lambda\approx 0{,}3994$ y $\lambda\approx 1001{,}6$ (\cref{eje:2-rosenbrock}). Cerca del mínimo, $\kappa\approx 2508$ y el descenso de gradiente solo es estable con $\eta<2/1001{,}6\approx 0{,}002$. El mínimo está en el fondo de un valle curvo, estrecho y casi plano a lo largo de su eje. Con $\eta=10^{-3}$ y desde $(-1{,}5;\,2)$, el método no baja de $f\approx 6{,}5\cdot 10^{-3}$ en 5000 iteraciones; es el banco de pruebas de la práctica del módulo.
\end{ejemplo}

\begin{observacion}[Borde de estabilidad]\label{obs:borde}
La condición $\eta<2/\lambda_{\max}$ también se observa en redes profundas, pero en sentido inverso. Cohen \emph{et al.}~\cite{cohen} observaron que, al entrenar con descenso de gradiente de paso fijo, la curvatura máxima $\lambda_{\max}$ del hessiano de la pérdida crece hasta situarse en torno a $2/\eta$ y permanece ahí; la pérdida deja de decrecer de forma monótona a corto plazo, pero sigue bajando a largo plazo. La red evoluciona hacia el límite de estabilidad del \cref{teo:gd-cuadratica}.
\end{observacion}

\begin{lecturas}
\item Boyd y Vandenberghe~\cite{bv}, §9.2 «Descent methods» y §9.3 «Gradient descent method».
\item MIT 18.065~\cite{mit18065}, clases~21 «Minimizing a Function Step by Step» y~22 «Gradient Descent: Downhill to a Minimum»; Strang~\cite{strang}, cap.~VI.
\item Zhang \emph{et al.}~\cite{d2l}, §12.3 «Gradient Descent».
\item Goodfellow, Bengio y Courville~\cite{gbc}, §4.3 «Gradient-Based Optimization» y §8.2.1 «Ill-Conditioning».
\end{lecturas}

\section*{Ejercicios del tema 2}
\addcontentsline{toc}{section}{Ejercicios del tema 2}

\begin{ejercicio}{1}{eje:2-escala}
Se multiplica la función objetivo por una constante $c>0$, es decir, se minimiza $cf$ en lugar de $f$. ¿Cómo cambian el minimizador, el paso máximo estable del descenso de gradiente y el número de condición? ¿Qué paso hay que usar con $cf$ para reproducir exactamente los iterados obtenidos con $f$ y paso $\eta$?
\end{ejercicio}

\begin{ejercicio}{2}{eje:2-1d}
Sea $f(x)=3x^{2}$, $x_0=2$ y $\eta=0{,}1$. Calcule $x_1$, $x_2$ y $x_3$. ¿Para qué valores de $\eta$ converge el método?
\end{ejercicio}

\begin{ejercicio}{2}{eje:2-2d}
Sea $f(x,y)=x^{2}+5y^{2}$ con punto inicial $(1,1)$ y $\eta=1/6$. Calcule tres iteraciones y el factor por el que se multiplica la norma del iterado en cada paso. Compruebe que coincide con el del \cref{teo:gd-cuadratica}.
\end{ejercicio}

\begin{ejercicio}{2}{eje:2-valle}
Sea $f(x,y)=\tfrac12(x^{2}+25y^{2})$. (a) Calcule $\kappa$, el paso óptimo $\eta^{\ast}$ y el factor $\rho$ del descenso de gradiente. (b) ¿Cuántas iteraciones hacen falta para reducir el error a $10^{-3}$? (c) Repita (b) con el factor del momentum óptimo, $(\sqrt\kappa-1)/(\sqrt\kappa+1)$. (d) Calcule los dos primeros iterados desde $(1,1)$ con $\eta^{\ast}$.
\fuente{adaptado de Zhang \emph{et al.}~\cite{d2l}, §12.3, ejercicio~3 (diseñar una función con coordenadas a escalas distintas en la que el descenso sea lento).}
\end{ejercicio}

\begin{ejercicio}{2}{eje:2-paso-largo}
Sea $F(x,y)=y^{2}+(y-x)^{2}$, cuyo mínimo es $F=0$ en $(0,0)$. (a) Calcule $\nabla F(1,1)$ y tres pasos de descenso de gradiente con paso $s=1/2$ desde $(1,1)$. ¿Qué ocurre con $F$? (b) Explique el resultado con los autovalores del hessiano. (c) Repita (a) con $s=1/4$.
\fuente{adaptado de MIT 18.065~\cite{mit18065}, problemas de la clase~25 (sección~VI.5 de~\cite{strang}), problema~2. Se usa la versión del enunciado cuyo mínimo es $(0,0)$.}
\end{ejercicio}

\begin{ejercicio}{2}{eje:2-rosenbrock}
Calcule los autovalores del hessiano de la función de Rosenbrock en $(1,1)$ (\cref{ej:rosenbrock}), su número de condición y el mayor paso estable del descenso de gradiente. Con $\eta=10^{-3}$, estime cuántas iteraciones hacen falta para reducir a $10^{-3}$ el error en la dirección de curvatura mínima.
\end{ejercicio}

\begin{ejercicio}{3}{eje:2-cuadratica}
Demuestre con detalle el \cref{teo:gd-cuadratica}, justificando en particular por qué el máximo de $\abs{1-\eta\lambda}$ sobre $[\mu,L]$ se alcanza en un extremo y por qué el valor óptimo de $\eta$ iguala los dos extremos.
\end{ejercicio}
